\chapter*{TÓM LƯỢC}
\addcontentsline{toc}{chapter}{TÓM LƯỢC}

Ô nhiễm không khí, đặc biệt bụi mịn $PM_{2.5}$, đang là vấn đề nghiêm trọng tại Việt Nam, ảnh hưởng trực tiếp đến sức khỏe cộng đồng. Việc dự báo chính xác nồng độ $PM_{2.5}$ trong ngắn hạn (1–24 giờ) đóng vai trò quan trọng trong hệ thống cảnh báo sớm và hỗ trợ ra quyết định. Đồ án này xây dựng và so sánh \textbf{6 mô hình} dự báo $PM_{2.5}$ multi-horizon ($T+1, T+3, T+6, T+12, T+24$ giờ), bao gồm: XGBoost, XLinear, ESTGCN, ST-XLinear, Ensemble và iTransformer. Dữ liệu được thu thập từ \textbf{12 trạm quan trắc} (6 miền Bắc, 6 miền Nam), kết hợp dữ liệu chất lượng không khí ($PM_{2.5}, PM_{10}, CO, NO_2, SO_2, O_3$) và dữ liệu khí tượng (nhiệt độ, độ ẩm, gió, lượng mưa, nhiệt độ đất). Một pipeline tiền xử lý hoàn chỉnh gồm 12 bước đã được thiết kế, bao gồm phát hiện lỗi cảm biến quang học (OPC), xử lý dữ liệu đóng băng, ngoại lai, và 6 đặc trưng kỹ thuật mới dựa trên kiến thức hóa học khí quyển (oxidation potential, humid sulfate risk, thermal stability...). Đặc biệt, đồ án đề xuất \textbf{Dynamic Graph Builder} — ma trận kề thay đổi theo hướng gió thực tế ($\alpha = 0.4858$), nhằm mô hình hóa sự lan truyền ô nhiễm giữa các trạm. Kết quả cho thấy \textbf{XGBoost đạt hiệu suất cao nhất} xuyên suốt mọi horizon dự báo, với $R^2 = 90.47\%$ ($T+1$) và $70.95\%$ ($T+24$), vượt trội so với các mô hình Deep Learning. Sai số RMSE trung bình chỉ dao động từ 5.16 $\mu g/m^3$ ($T+1$) đến 11.52 $\mu g/m^3$ ($T+24$). Cụm trạm miền Nam cho kết quả tốt hơn miền Bắc ($R^2 = 98.81\%$ vs $82.12\%$ tại $T+1$), phản ánh mức độ biến động của hạt $PM_{2.5}$ thấp hơn nhờ đặc thù khí hậu rửa trôi. Toàn bộ hệ thống đã được tích hợp thành ứng dụng Web Streamlit, cho phép giám sát và dự báo trực tuyến $PM_{2.5}$ tại 12 trạm, sẵn sàng phục vụ mục đích cảnh báo dân sinh.